\begin{ouverture}[calcul]{5}{Fractions, quotients et comparaisons}{Partie 1 · Nombres et calculs}
{Un nombre · Plusieurs écritures · Dénominateur commun}
\objectif{Comprendre une fraction comme nombre quotient.}
\objectif{Passer entre des écritures équivalentes.}
\objectif{Comparer et placer des fractions.}
\prerequis{Multiples, divisions et droites graduées.}
\end{ouverture}
\pagestyle{otMain}\essentiel\fontsize{11}{15}\selectfont
\phantomsection\label{cours-5-1}
\essSection{1}{Un quotient est un nombre}
\begin{definition}Pour $b\ne0$, le quotient $a/b$ est le nombre qui, multiplié par $b$, donne $a$. Une fraction est un quotient de deux entiers, avec un dénominateur non nul. Le numérateur compte des parts ; le dénominateur indique le partage de l’unité.\end{definition}
\label{fig-cours-5-1-cours-fractions}\figureGeo{cours-fractions}

\begin{propriete}Multiplier ou diviser numérateur et dénominateur par un même nombre non nul conserve la valeur : $\frac34=\frac68=\frac{15}{20}$. Une fraction peut être plus grande que 1 : $\frac74=1+\frac34$.\end{propriete}
\begin{exemple}$1{,}25=\frac{125}{100}=\frac54=1+\frac14=125\,\%$. Ce sont différentes écritures du même nombre. Un quotient comme $1/3$ ne possède pas d’écriture décimale finie ; $0{,}33$ n’en est qu’une approximation.\end{exemple}
\clearpage
\phantomsection\label{cours-5-2}
\essSection{2}{Choisir une méthode de comparaison}
\begin{propriete}À dénominateur positif identique, la plus grande fraction a le plus grand numérateur. À numérateur positif identique, la fraction ayant le plus petit dénominateur est la plus grande.\end{propriete}
\begin{savoirfaire}[Un même partage]Pour comparer $\frac58$ et $\frac7{12}$, choisir un dénominateur commun : $\frac{15}{24}>\frac{14}{24}$. On peut aussi comparer à 1 ou à $1/2$ quand cela suffit.\end{savoirfaire}
\begin{exemple}Pour placer $11/4$, effectuer $11=4\times2+3$ : $\frac{11}4=2+\frac34$. Partager l’intervalle de 2 à 3 en quatre et compter trois parts.\end{exemple}
\begin{attention}Le dénominateur ne se lit pas comme une partie indépendante du nombre. $\frac25<\frac23$ malgré $5>3$ : cinquièmes et tiers ne sont pas les mêmes parts. Ne pas comparer uniquement les chiffres.\end{attention}

\clearpage\rubrique{Découverte}
\begin{blocLivre}{ex-5-1}\begin{exercice}[1]{Compléter un produit}
Quel nombre complète $7\times\,?=5$ ?
\end{exercice}\end{blocLivre}
\begin{blocLivre}{ex-5-2}\begin{exercice}[1]{Fractions égales}
Compléter $\frac35=\frac{\,?\,}{20}$ et $\frac47=\frac{12}{\,?\,}$.
\end{exercice}\end{blocLivre}
\begin{blocLivre}{ex-5-3}\begin{exercice}[1]{Écriture décimale}
Écrire $\frac12$, $\frac34$ et $\frac52$ sous forme décimale.
\end{exercice}\end{blocLivre}
\begin{blocLivre}{ex-5-4}\begin{exercice}[1]{Même dénominateur}
Comparer $7/11$ et $9/11$.
\end{exercice}\end{blocLivre}
\begin{blocLivre}{ex-5-5}\begin{exercice}[1]{Même numérateur}
Comparer $5/8$ et $5/13$.
\end{exercice}\end{blocLivre}

\rubrique{Entraînement}
\begin{blocLivre}{ex-5-6}\begin{exercice}[2]{Un dénominateur commun}
Comparer $5/6$ et $7/9$.
\end{exercice}\end{blocLivre}
\begin{blocLivre}{ex-5-7}\begin{exercice}[2]{Dépasser l’unité}
Écrire $23/6$ comme somme d’un entier et d’une fraction inférieure à 1.
\end{exercice}\end{blocLivre}
\begin{blocLivre}{ex-5-8}\begin{exercice}[2]{Placer des fractions}
Placer $1/2$, $5/4$ et $7/4$ sur la droite graduée en quarts.
\label{fig-ex-5-8}\figureGeo{frac-8}
\end{exercice}\end{blocLivre}
\begin{blocLivre}{ex-5-9}\begin{exercice}[2]{Des formes équivalentes}
Écrire 0,6 comme fraction de dénominateur 10, fraction de dénominateur 5 et pourcentage.
\end{exercice}\end{blocLivre}
\begin{blocLivre}{ex-5-10}\begin{exercice}[2]{Exact ou approché}
$2/3=0{,}67$ est-il exact ? Justifier sans calculatrice.
\end{exercice}\end{blocLivre}

\rubrique{Problèmes}
\begin{blocLivre}{ex-5-11}\begin{exercice}[3]{Classer trois parts}
Ranger $3/4$, $5/6$ et $7/12$ dans l’ordre croissant.
\end{exercice}\end{blocLivre}
\begin{blocLivre}{ex-5-12}\begin{exercice}[3]{Comparer sans tout calculer}
Comparer $11/10$ et $9/11$, puis $4/9$ et $5/8$ en utilisant 1 ou $1/2$.
\end{exercice}\end{blocLivre}
\begin{blocLivre}{ex-5-13}\begin{exercice}[3]{Trois recettes}
Dans A, il y a 3 doses de sirop pour 8 doses de boisson ; dans B, 5 pour 12 ; dans C, 2 pour 5. Quelle boisson a la plus grande proportion de sirop ?
\end{exercice}\end{blocLivre}
\begin{blocLivre}{ex-5-14}\begin{exercice}[3]{Un encadrement fin}
Trouver tous les entiers $a$ tels que $\frac23<\frac a{12}<\frac56$.
\end{exercice}\end{blocLivre}
\begin{blocLivre}{ex-5-15}\begin{exercice}[3]{Deux affichages, une valeur}
Une jauge affiche 1,5 unité, une autre $6/4$, une troisième 150 \%. Expriment-elles le même nombre ? Que faut-il vérifier si on compare des remplissages ?
\end{exercice}\end{blocLivre}

\clearpage
